% !TEX root = ../main.tex
\section{SQL Templates, Configuration, and Extended Notes}
\label{ch:appendix-sql}

To aid reproduction, this appendix shows the three SQL patterns on which every graph of the thesis rests, the configuration keys, and notes on complexity and on the conventions of the result tables. The listings are abridged; the full files are in \texttt{sql/} of the experiments folder (Appendix~\ref{ch:appendix-eval}). Placeholders in double underscores, such as \texttt{\_\_DEST\_\_}, are replaced by \texttt{scripts/run\_bq.py} before a query runs, and names that start with \texttt{@} are query parameters.

\dissertationSubheading[sec:sql-materialize]{Log materialization pattern}

Scan~A and scan~B run the same query, \texttt{sql/01\_materialize\_logs.sql}, over different windows. It reads the public log table once, tags each log by kind and keeps only the four kinds the study needs, in a project table partitioned by month and clustered by kind. With the GMX and Aave branches shortened, it reads:

\begin{verbatim}
CREATE TABLE `__DEST__`
PARTITION BY TIMESTAMP_TRUNC(block_timestamp, MONTH)
CLUSTER BY kind
AS
WITH src AS (
  SELECT block_timestamp, block_number, log_index,
         LOWER(address) AS address, topics, data
  FROM `bigquery-public-data.goog_blockchain_arbitrum_one_us.logs`
  WHERE block_timestamp >= TIMESTAMP(@start_ts)
    AND block_timestamp < TIMESTAMP(@end_ts)
),
tagged AS (
  SELECT *,
    CASE
      WHEN ARRAY_LENGTH(topics) = 3 AND topics[OFFSET(0)] =
        '0x8c5be1e5ebec7d5bd14f71427d1e84f3dd0314c0f7b2291e5b200ac8c7c3b925'
        THEN 'A'                                  -- ERC-20 Approval
      WHEN ARRAY_LENGTH(topics) = 3 AND topics[OFFSET(0)] =
        '0xddf252ad1be2c89b69c2b068fc378daa952ba7f163c4a11628f55a4df523b3ef'
        THEN 'T'                                  -- ERC-20 Transfer
      WHEN address = '0xc8ee91a54287db53897056e12d9819156d3822fb'
           AND ... THEN 'G'                       -- GMX V2 PositionDecrease
      WHEN address = '0x794a61358d6845594f94dc1db02a252b5b4814ad'
           AND ... THEN 'V'                       -- Aave V3 pool events
    END AS kind
  FROM src
)
SELECT
  block_timestamp, block_number, log_index, kind,
  FROM_HEX(SUBSTR(address, 3)) AS emitter,
  IF(kind IN ('A', 'T'), FROM_HEX(SUBSTR(topics[OFFSET(1)], 27)), NULL) AS a,
  IF(kind IN ('A', 'T'), FROM_HEX(SUBSTR(topics[OFFSET(2)], 27)), NULL) AS b,
  IF(kind IN ('A', 'T'), SAFE.FROM_HEX(SUBSTR(data, 3)), NULL) AS val,
  IF(kind IN ('G', 'V'), topics, NULL) AS raw_topics,
  IF(kind IN ('G', 'V'), data, NULL) AS raw_data
FROM tagged
WHERE kind IS NOT NULL
\end{verbatim}

For an approval, \texttt{a} is the owner and \texttt{b} the spender; for a transfer, \texttt{a} is the sender and \texttt{b} the recipient. In both, \texttt{emitter} is the token contract and \texttt{val} the 32-byte amount. The two topic-0 hashes are the keccak256 hashes of \texttt{Approval(address,address,uint256)} and \texttt{Transfer(address,address,uint256)} \parencite{vogelsteller2015erc20}. Requiring exactly three topics keeps the ERC-20 logs and leaves out the ERC-721 logs, which share both hashes but index the token id as a fourth topic. Addresses are stored as 20-byte values, not as hexadecimal strings.

\dissertationSubheading[sec:sql-ego]{Ego-network extraction pattern}

\texttt{sql/03\_ego\_events.sql} reads the materialized table of one window and keeps the approvals and transfers with a cohort contract at one end. The table \texttt{contract\_rep.cohort} holds the 20-byte addresses of the matched cohort, loaded by \texttt{scripts/build\_cohort.py}. Abridged:

\begin{verbatim}
CREATE TEMP FUNCTION u256(v BYTES) AS (
  IFNULL((SELECT SUM(c * POW(256.0, LENGTH(v) - 1 - o))
          FROM UNNEST(TO_CODE_POINTS(v)) AS c WITH OFFSET o), 0.0)
);

CREATE OR REPLACE TABLE `dissertation-bq.contract_rep.approvals___SUFFIX__`
PARTITION BY TIMESTAMP_TRUNC(block_timestamp, MONTH)
AS
SELECT
  l.block_timestamp, l.block_number, l.log_index,
  l.emitter AS token_address,
  l.a AS owner,
  l.b AS spender,
  u256(l.val) AS value
FROM `dissertation-bq.contract_rep.__SRC__` AS l
WHERE l.kind = 'A'
  AND l.block_timestamp >= TIMESTAMP(@start_ts)
  AND l.block_timestamp < TIMESTAMP(@end_ts)
  AND (l.a IN (SELECT address FROM `dissertation-bq.contract_rep.cohort`)
       OR l.b IN (SELECT address FROM `dissertation-bq.contract_rep.cohort`));

-- transfers___SUFFIX__: the same with kind = 'T', from_address = a and
-- to_address = b, and zero amounts dropped:
--   AND LTRIM(TO_HEX(l.val), '0') != ''
\end{verbatim}

The function \texttt{u256} reads the 32-byte amount as a double in raw base units; no token decimals are applied anywhere in the pipeline. Approvals of zero are kept, because an approval of zero deletes an allowance, while zero-value transfers are dropped. The suffix is \texttt{ego} for the observation window and \texttt{w1ego} for July--September~2026.

\dissertationSubheading[sec:sql-pairs]{Latest-allowance and pair-table pattern}

\texttt{sql/04\_anchor\_pairs.sql} turns the ego events up to one anchor into the edge lists of both graphs. Abridged:

\begin{verbatim}
CREATE TEMP FUNCTION sig(age_days FLOAT64) AS
  (1.0 / (1.0 + EXP(0.01 * (age_days - 180.0))));
CREATE TEMP FUNCTION vt(x FLOAT64) AS
  (IF(x >= 40.0, 1.0, 2.0 / (1.0 + EXP(-x)) - 1.0));

CREATE OR REPLACE TABLE `dissertation-bq.contract_rep.allow_triples___TAG__` AS
SELECT token_address, owner, spender, value, block_timestamp,
       TIMESTAMP_DIFF(TIMESTAMP(@anchor_ts), block_timestamp, MICROSECOND)
         / 8.64e10 AS age_days
FROM (
  SELECT *, ROW_NUMBER() OVER (
           PARTITION BY token_address, owner, spender
           ORDER BY block_number DESC, log_index DESC) AS rn
  FROM `dissertation-bq.contract_rep.approvals_ego`
  WHERE block_timestamp >= TIMESTAMP(@start_ts)
    AND block_timestamp <= TIMESTAMP(@anchor_ts)
)
WHERE rn = 1 AND value > 0;

CREATE OR REPLACE TABLE `dissertation-bq.contract_rep.allow_pairs___TAG__` AS
SELECT owner, spender,
       COUNT(*) AS n_tokens,
       SUM(sig(age_days) * vt(value)) AS w_er,   -- EndorseRank
       SUM(value) AS w_raw,                      -- C-PR allowance layer
       MAX(sig(age_days)) AS max_sig,
       MAX(block_timestamp) AS last_ts
FROM `dissertation-bq.contract_rep.allow_triples___TAG__`
GROUP BY owner, spender;

-- transfer_pairs___TAG__: per (from_address, to_address) over transfers_ego
-- up to the anchor, with
--   SUM(sig(age_days) * vt(value)) AS w_awp    -- AWP
--   SUM(value * sig(age_days))     AS w_raw    -- C-PR transfer layer
--   MAX(sig(age_days))             AS max_sig  -- AWP activeness
\end{verbatim}

The function \texttt{sig} is the logistic decay of Equation~\eqref{eq:logistic-decay} with $k=0.01$ and $t_0=180$ days, and \texttt{vt} the value transform of Equation~\eqref{eq:value-transform} with $b=1$, set to one from $z=40$ on, where it equals one in double precision. Among the approvals of one token, owner and spender up to the anchor, the latest is the one with the largest block number and, within the block, the largest log index. It is kept only when its value is positive, so a revocation removes that token from the pair. Each pair row carries the weights of every method, so all scores at one anchor come from the same rows. The tag is \texttt{tobs} for $T_{\text{obs}}$ and \texttt{t1} for the freeze of W0.

\dissertationSubheading[sec:config-keys]{Configuration keys}

All of the pipeline's settings are kept in one versioned file, \texttt{config/contract\_reputation.yaml}, which the scripts read at start-up. Its groups of keys are:

\begin{itemize}
\item \texttt{chain\_id}, \texttt{chain\_name} and \texttt{bigquery}: chain 42161 (Arbitrum One); the project, dataset, location and public log table; the budget in USD; and the windows of scan~A and scan~B;
\item \texttt{period}, \texttt{holdout} and \texttt{registered\_window}: the observation window with its anchor $T_{\text{obs}}$, the freeze and label window of W0, and the freeze, label window, status and margin $\delta$ of W1;
\item \texttt{events}, \texttt{gmx\_arbitrum} and \texttt{aave\_arbitrum}: the topic-0 hashes of \texttt{Approval} and \texttt{Transfer}; the GMX~V2 event emitter, the \texttt{EventLog1} topic, the \texttt{PositionDecrease} hash, the liquidation handler, the order type of a liquidation and the minimum of three closes; and the Aave~V3 pool with its event topics;
\item \texttt{accounts} and \texttt{cohort}: the RPC endpoint and the delegation prefix \texttt{0xef0100} of the code check, and the minimum of three distinct owners;
\item \texttt{reputation}: damping, tolerance, iteration cap, decay $k$ and $t_0$, value parameter $b$, the primary and sensitivity values of $\lambda$, and the two parts of rule~A;
\item \texttt{alignment}, \texttt{benchmark} and \texttt{robustness}: the bootstrap resamples and seed; the warm-up and timed runs and the scaling seed; the damping values and the number of top tokens;
\item \texttt{storage}: the largest size of a data file part, 90\,MB.
\end{itemize}

Changing the damping factor, the decay or the value parameter requires running the evaluation again, the damping, top-token and sample-definition checks included, and exporting the tables and figures again from the same summaries, so that every reported number matches its machine-readable source. \texttt{scripts/extract\_w1.py} records the SHA-256 hash of this file at the start of scan~B (Appendix~\ref{sec:map-artifacts}).

\dissertationSubheading[sec:complexity-notes]{Complexity notes}

Write $m=|E|$. In BigQuery the latest allowances come from a window function over the $R_a$ approval rows, which sorts within each token, owner and spender, in $O(R_a\log R_a)$ time at most; the transfer pairs are one aggregation over the $R_t$ transfer rows, in $O(R_t)$. Locally, PageRank over $I$ iterations costs $O(m+mI)$, the first term for building the sparse operator and the second for one sparse product per iteration (Appendix~\ref{sec:pr-sparsity}). C-PR walks the union of both layers, so one of its iterations costs as much as one on each layer, and S-PR solves the allowance layer first and then the transfer layer. The edge lists and score vectors take most of the memory, so the peak memory of a solve grows with $m$; in Table~\ref{tab:benchmark-runtime} the peak-memory ratio of AWP to EndorseRank is $2.33$ ($2{,}995.9$ against $1{,}285.0$~MB), at an edge ratio of $2.68$.

\dissertationSubheading[sec:interpretation-checklist]{Conventions for the result tables}

The tables of Chapter~\ref{ch:results} and of this appendix follow these conventions:

\begin{enumerate}
\item Every $\tau$ is Kendall's $\tau_b$ on raw scores and proxies, reported with its 95\% paired-bootstrap interval: $400$ resamples over contracts with seed 42, and $2{,}000$ for the trader comparisons on outcomes built from neither edge type, whose decision contrasts P and Q also carry 97.5\% intervals. Differences between coefficients come from the paired contrast tables, not from subtracting point estimates.
\item A cohort contract missing from a method's graph scores zero under that method. Keeping it as an isolated node instead is the sensitivity analysis of Table~\ref{tab:tie-sensitivity}.
\item Each contrast table says which of its contrasts were fixed in advance, by the plan (rule~A) or by the registration of W1 (rule~B and the comparison on outcomes built from neither edge type), and which were computed after the labels were known.
\item A runtime is the mean wall-clock time of five timed solver calls after one untimed warm-up. A call turns edge lists built in advance into scores: node indexing, sparse-matrix assembly, row normalization, restart vector and power iteration, without extraction or preprocessing. Every runtime appears with $|V|$ and $|E|$ of its solver graph, and the same graph under the same settings shows the same measurement in every table (Section~\ref{sec:computational-benchmark}).
\item The sample-definition stages (Table~\ref{tab:robustness-sample-size}) show how the statistic moves with the evaluation set; they are not robustness evidence, and their values are not compared with the full-cohort tables.
\end{enumerate}

\dissertationSubheading[sec:limitations-reprise]{Limitations}

The queries above were run on Arbitrum One only. Their results depend on the public log table as it stood on 9~October~2026, when scan~A and scan~B ran. The ego filter keeps only events with a cohort contract at one end, so the graphs contain no edge between two addresses outside the cohort. All reported numbers come from the machine-readable summaries of the matched cohort, the spender and trader cohorts of W0 and W1, the scaling stages and the Sybil model; the registered results of July--September come from the summary of that run (Section~\ref{sub:fresh-results}).
